\documentclass[11pt]{article}
\usepackage[T1]{fontenc}
\usepackage{lmodern,amsmath,amssymb,amsthm,mathtools,microtype}
\usepackage[margin=1in]{geometry}
\usepackage{xcolor}
\usepackage[colorlinks=true,allcolors=blue!45!black]{hyperref}
\hypersetup{pdftitle={Absolute error of fractional von Mangoldt approximants on shifted primes}}
\newtheorem{theorem}{Theorem}[section]
\newtheorem{lemma}[theorem]{Lemma}
\newtheorem{proposition}[theorem]{Proposition}
\newtheorem{corollary}[theorem]{Corollary}
\theoremstyle{remark}
\newtheorem{remark}[theorem]{Remark}
\newcommand{\PP}{\mathbb P}
\newcommand{\E}{\mathbb E}
\newcommand{\one}{\mathbf 1}
\newcommand{\rad}{\operatorname{rad}}
\newcommand{\Li}{\operatorname{Li}}
\newcommand{\TV}{d_{\mathrm{TV}}}
\newcommand{\err}{\mathcal E_X}
\newcommand{\Mcal}{\mathcal M_X}
\allowdisplaybreaks[1]
\raggedbottom
\setlength\emergencystretch{1.5em}
\title{\bfseries Absolute error of fractional von Mangoldt\\approximants on shifted primes}
\author{\normalsize DEEPAK KUMAR (research coordinator)\\[-2pt]\small\href{mailto:admin@cosmosgeek.org}{admin@cosmosgeek.org}}
\date{\small Public working paper\quad 25 November 2026}
\begin{document}
\maketitle

\begin{abstract}
For real $\alpha>0$ put
$Q_\alpha(m)=\sum_{d\mid m}\mu(d)(\log(m/d)/\log m)^\alpha$.
These weights equal one on primes and $Q_1(m)=\Lambda(m)/\log m$.
On the shifted-prime row $m=p+2$, we evaluate the entire composite
absolute error at orders $1\pm c/\log\log X$: in the normalization
$2C_2X/\log^2X$, its limits are $1-e^{-c}$ and $e^c-1$.
A fixed binomial combination of orders strictly greater than one has
absolute error $(1-e^{-1})^k$ in the same normalization.
The arithmetic input is Ford's total-variation model for small prime
valuations, together with the classical upper-bound sieve. We give the
uniform positive envelope needed to integrate the moving-order Mellin
formula, including repeated prime factors. The estimates do not supply
a positive lower bound for the twin-prime count.
\end{abstract}

\section{Statements and arithmetic inputs}
Let $X\to\infty$, $L=\log X$, $H=\log L$, and
\[
 C_2=\prod_{\ell>2}\left(1-\frac1{(\ell-1)^2}\right),\qquad
 \mathcal P_X=\{p\in\PP:X<p\le2X\}.
\]
Letters $p,\ell$ denote primes. For $m>1$ and real $\alpha>0$ define
\begin{equation}\label{eq:Q}
 Q_\alpha(m)=\sum_{d\mid m}\mu(d)
       \left(\frac{\log(m/d)}{\log m}\right)^\alpha,
 \qquad 0^\alpha=0.
\end{equation}
For a function $F$ with $F(p)=1$ at primes, write
\begin{equation}\label{eq:error}
 \err(F)=\frac{L^2}{2C_2X}\sum_{p\in\mathcal P_X}
        |F(p+2)-\one_{\PP}(p+2)|.
\end{equation}
Thus proper prime powers are included in this definition of error.

\begin{theorem}\label{th:critical}
For every fixed $c>0$,
\[
 \err(Q_{1+c/H})\longrightarrow1-e^{-c},\qquad
 \err(Q_{1-c/H})\longrightarrow e^c-1.
\]
The second order is positive for all sufficiently large $X$.
For every fixed positive integer $k$, the function
\begin{equation}\label{eq:Rk}
 R_{k,X}(m)=\sum_{j=1}^k(-1)^{j+1}\binom{k}{j}Q_{1+j/H}(m)
\end{equation}
equals one on primes and satisfies
\[
       \err(R_{k,X})\longrightarrow(1-e^{-1})^k.
\]
All parameters $c$ and $k$ are fixed before $X$ tends to infinity.
\end{theorem}

For comparison, set
\[
 \tau_X=\frac{L^2}{2C_2X}
       |\{p\in\mathcal P_X:p+2\in\PP\}|,
 \qquad
 \Mcal(F)=\frac{L^2}{2C_2X}\sum_{p\in\mathcal P_X}F(p+2).
\]
One always has $|\Mcal(F)-\tau_X|\le\err(F)$. Theorem~\ref{th:critical}
evaluates the right side. It does not evaluate the signed quantity
$\Mcal(F)$ or show that it exceeds this error. Nor does it justify
choosing $k=k(X)$ without a further uniform estimate.

The functions $Q_k$ at positive integer orders are normalized generalized
von Mangoldt functions. Higher moments and their distinction from prime
detection are classical; see Ford's discussion of the asymptotic sieve
\cite{ford-asymptotic}. The distribution input here is much more directly
Ford's valuation model \cite[Theorem 1]{ford-model}. The contribution of
this note is the explicit moving fractional-order absolute-error law
and its binomial consequence, as applications of that model. We assert
no new level of distribution or prime-factor independence theorem, and
make no global priority claim for this formulation.

We record the precise form of Ford's input. Let $W_2=0$ and let the
$W_\ell$, for odd primes $\ell$, be independent with
\begin{equation}\label{eq:ford-law}
 \mathbb P(W_\ell=0)=1-\frac1{\ell-1},\qquad
 \mathbb P(W_\ell=a)=\ell^{-a}\quad(a\ge1).
\end{equation}
For a uniformly chosen prime $p\le x$, with finitely many small primes
omitted, Ford's theorem gives
\begin{equation}\label{eq:ford}
 \TV\big((v_\ell(p+2))_{\ell\le y},(W_\ell)_{\ell\le y}\big)
 \ll_{A,\theta} e^{-\theta u\log u}+\log^{-A}x,
 \qquad u=\frac{\log x}{\log y},
\end{equation}
for $2\le y\le x$, every fixed $A>0$ and any fixed $0<\theta<1/2$.
Indeed the theorem assumes its hypothesis $Z(\gamma)$ and
$\theta<\gamma$; $Z(1/2)$ follows from Bombieri--Vinogradov, as stated
in \cite{ford-model}. The shift is the fixed integer $2$.
Subtracting the cumulative estimates at $2X$ and $X$ gives the same
bound, up to a constant and the smaller of the two $u$'s, for
$p\in\mathcal P_X$. The ratio of the cumulative prime counts to
$|\mathcal P_X|$ is bounded by the prime number theorem.
In particular, \eqref{eq:ford} applies to every $[0,1]$-valued function
of the indicated valuations, even when that function depends on $X$.

\section{A uniform nonnegative envelope}
For $w>0$ define the nonnegative multiplicative weights
\begin{equation}\label{eq:DH}
 D_w(m)=\prod_{\ell\mid m}(1-\ell^{-w}),\qquad
 B_w(m)=\left(\frac{\rad m}{m}\right)^wD_w(m).
\end{equation}
Both weights lie in $[0,1]$. Let
\[
 \mathcal C_X=\{p\in\mathcal P_X:\omega(p+2)\ge2\},
\]
where $\omega$ counts distinct prime divisors. For every $t>0$, put
\begin{equation}\label{eq:AX}
 A_X(t)=\frac{L^2}{2C_2X}\sum_{p\in\mathcal C_X}
                  B_{t/\log(p+2)}(p+2).
\end{equation}

\begin{lemma}[Randomized upper sieve]\label{lem:sieve}
Fix $a=1/100$ and $z=X^a$. For arbitrary numbers
$g_\ell\in[0,1]$ on the odd primes $\ell\le z$,
\begin{equation}\label{eq:random-upper}
 \sum_{p\in\mathcal P_X}\prod_{\substack{\ell\mid p+2\\2<\ell\le z}}g_\ell
 \ll \frac{X}{L^2}
       \exp\left(\sum_{2<\ell\le z}\frac{g_\ell}{\ell-2}\right)
       +X^{1/2}.
\end{equation}
The implicit constant is independent of $X$ and of all the $g_\ell$.
\end{lemma}

\begin{proof}
Independently put each odd prime $\ell\le z$ in a set $\mathcal S$ with
probability $1-g_\ell$. The product in \eqref{eq:random-upper} is the
probability that no member of $\mathcal S$ divides $p+2$.
For a fixed $\mathcal S$, majorize the prime condition by excluding
$0\pmod\ell$ from $n$ for every $\ell\le z$, and also exclude
$-2\pmod\ell$ when $\ell\in\mathcal S$. At $2$ there is one excluded
residue; at an odd prime there are
$\rho_\ell=1+\one_{\mathcal S}(\ell)$.

The upper-bound fundamental lemma, in the form of
\cite[Theorem 19.1 and its deduction of Theorem 18.11]{kouk}, gives
upper sieve weights supported on squarefree $d\le D=X^{1/4}$,
with absolute value at most one, and main term at most a fixed
constant times $X\prod_{\ell\le z}(1-\rho_\ell/\ell)$.
Here the sieve clock is $25$, and the dimension-two product condition
is uniform because $\rho_2=1$ and $1\le\rho_\ell\le2$ at odd primes.
The latter condition follows directly from Mertens' estimate applied
to the larger factors $(1-2/\ell)^{-1}$, with the prime $2$ separated.
For completeness, CRT gives $\rho(d)$ forbidden-residue combinations
and discrepancy $O(\rho(d))$ on the interval $(X,2X]$. Therefore the
remainder is at most
\[
 \sum_{d\le D}\mu^2(d)2^{\omega(d)}\ll D\log(2D)\ll X^{1/2},
\]
uniformly in $\mathcal S$. This is an integer-sequence sieve, not a
distribution assumption about primes.

Taking expectations, its local main factors become
\[
 \frac12\prod_{2<\ell\le z}
       \left(1-\frac2\ell+\frac{g_\ell}{\ell}\right)
 =\frac12\prod_{2<\ell\le z}\left(1-\frac2\ell\right)
       \prod_{2<\ell\le z}\left(1+\frac{g_\ell}{\ell-2}\right).
\]
The first product is $\ll1/\log^2z\ll1/L^2$, and the second is bounded
by the exponential in \eqref{eq:random-upper}.
\end{proof}

\begin{proposition}[Uniform envelope]\label{prop:envelope}
For all sufficiently large $X$, with one absolute implicit constant,
\begin{equation}\label{eq:envelope}
 A_X(t)\ll
 \begin{cases}
 t^2,&0<t\le1,\\
 t,&1\le t\le L,\\
 L,&t\ge L.
 \end{cases}
\end{equation}
Repeated prime factors are included throughout.
\end{proposition}

\begin{proof}
Suppose first that $0<t\le1$ and $m=p+2$, $p\in\mathcal C_X$.
Writing $r=\omega(m)\ge2$, the inequality $1-e^{-u}\le u$ gives
\[
 B_{t/\log m}(m)
 \le t^r\prod_{\ell\mid m}\frac{\log\ell}{\log m}
 \le t^2\prod_{\substack{\ell\mid m\\2<\ell\le z}}
                     \frac{\log\ell}{L}.
\]
In the last step all deleted factors are at most one, and $\log m>L$.
Apply Lemma~\ref{lem:sieve} with $g_\ell=\log\ell/L$.
The exponent is bounded, since
$L^{-1}\sum_{\ell\le z}\log\ell/\ell=O(1)$ and the replacement of
$1/\ell$ by $1/(\ell-2)$ costs $O(1)$.
Its remainder, after normalization and multiplication by $t^2$, is
$O(t^2L^2X^{-1/2})$. This proves the first bound.

For $1\le t\le L$ set $w=t/L$. Since $\log m\ge L$,
\[
 B_{t/\log m}(m)\le
 \prod_{\substack{\ell\mid m\\2<\ell\le z}}(1-\ell^{-w}).
\]
Use $g_\ell=1-\ell^{-w}$. Uniformly for $1\le t\le L$,
\begin{equation}\label{eq:harmonic-envelope}
 \sum_{2<\ell\le z}\frac{1-\ell^{-t/L}}{\ell-2}
                   \le\log t+O(1).
\end{equation}
To verify this, replace the denominator by $\ell$ at bounded cost
and split at $y_0=\exp(L/t)$. Below $y_0$ use
$1-\ell^{-t/L}\le(t/L)\log\ell$, whose reciprocal sum is $O(1)$.
If $y_0<z$, the remaining reciprocal-prime sum is at most
$\log(\log z/\log y_0)+O(1)\le\log t+O(1)$ by Mertens.
If $y_0\ge z$, the whole sum is $O(1)$.
The second bound follows from the randomized sieve; the normalized
remainder $O(L^2X^{-1/2})$ is absorbed by $t$.
Finally $B_w\le1$ and $|\mathcal P_X|\ll X/L$ give the third bound.
\end{proof}

\section{The positive interior asymptotic}
\begin{proposition}\label{prop:interior}
For every fixed $0<v<1$,
\begin{equation}\label{eq:interior}
                    \frac{A_X(L^v)}{L^v}\longrightarrow1.
\end{equation}
\end{proposition}

\begin{proof}
Put $t=L^v$, $w=t/L$, and $y=\exp(L/\sqrt t)$.
Then $w\to0$, $w\log y=\sqrt t$, and
$u=\log X/\log y=\sqrt t$ in Ford's theorem.
For a prime $\ell>2$ the model expectations of the local factors of
$D_w$ and $B_w$, respectively, are
\begin{align}
 d_\ell(w)&=1-\frac{\ell^{-w}}{\ell-1},\label{eq:local-d}\\
 b_\ell(w)&=1-\frac1{\ell-1}
       +\frac{1-\ell^{-w}}{\ell(1-\ell^{-1-w})}.
       \label{eq:local-b}
\end{align}
The difference satisfies
\begin{align*}
 0\le d_\ell(w)-b_\ell(w)
 &=\sum_{a\ge2}\ell^{-a}(1-\ell^{-w})
                         (1-\ell^{-(a-1)w})\\
 &\le\frac{w\log\ell\,(1-\ell^{-w})}{(\ell-1)^2}
 \le\frac{w^2(\log\ell)^2}{(\ell-1)^2}.
\end{align*}
Since $d_\ell(w)\ge1/2$, it follows uniformly in $y$ that
\begin{equation}\label{eq:model-powers}
       \prod_{2<\ell\le y}b_\ell(w)
          =\bigl(1+O(w^2)\bigr)\prod_{2<\ell\le y}d_\ell(w).
\end{equation}
The sign of the $O$ term here is immaterial.

The full undamped Euler product has the asymptotic
\begin{equation}\label{eq:euler-main}
 \prod_{\ell>2}d_\ell(w)\sim2C_2w\qquad(w\downarrow0).
\end{equation}
Indeed divide it by the odd-prime part of $\zeta(1+w)^{-1}$.
The quotient converges uniformly near $w=0$ by $O(\ell^{-2})$ local
errors, and its value at zero is $C_2$. The omitted prime $2$ supplies
$(1-2^{-1-w})^{-1}\to2$, while $\zeta(1+w)^{-1}\sim w$.
Also
\[
 \sum_{\ell>y}\frac{\ell^{-w}}{\ell-1}
 \ll\int_y^\infty\frac{x^{-1-w}}{\log x}\,dx
 \ll\frac{e^{-w\log y}}{w\log y}=o(1),
\]
by an elementary prime-counting upper bound and partial summation.
Consequently the product truncated at $y$ has the same asymptotic.

Both truncated multiplicative weights are $[0,1]$-valued functions of
the valuations through $y$. Apply \eqref{eq:ford}, with for example
$\theta=1/4$ and $A=4$. Its error is $o(w)$, so, averaging over
$p\in\mathcal P_X$,
\begin{equation}\label{eq:truncated-mean}
 \E_p D_w^{(\le y)}(p+2)\sim2C_2w,
 \qquad
 \E_p B_w^{(\le y)}(p+2)\sim2C_2w.
\end{equation}

We justify completing both weights. There are at most
$\log(2X+2)/\log y=O(\sqrt t)$ distinct prime divisors greater than
$y$ in any $m\le2X+2$. Hence
\[
 1-O(\sqrt t\,e^{-\sqrt t})
 \le\prod_{\substack{\ell\mid m\\\ell>y}}(1-\ell^{-w})\le1.
\]
For $B_w$ this also suffices unless a prime $\ell>y$ has
$\ell^2\mid m$. The number of such integers $m\le2X+2$ is
\[
 O\left(\sum_{y<\ell\le\sqrt{2X+2}}
                    (X/\ell^2+1)\right)
           =O(X/y+\sqrt X).
\]
After division by $|\mathcal P_X|\asymp X/L$ this is $o(w)$.
Thus \eqref{eq:truncated-mean} holds for the full $D_w,B_w$.
In particular $\E_p(D_w-B_w)(p+2)=o(w)$; this also follows directly
from \eqref{eq:model-powers} and the same completion.

The physical denominator causes no unproved uniformity assumption.
Write $w_m=t/\log m$ and $w_-=t/\log(2X+2)$, so
$w_-\le w_m\le w$ and $w_-/w\to1$. The preceding proof applies to
$w_-$ with the same $y$, since $w_-\log y\sim\sqrt t$.
Monotonicity of $D_w$ gives
\[
 \E_pD_{w_m}(p+2)\sim2C_2w.
\]
For each $m$, the difference
\[
 D_s(m)-B_s(m)=D_s(m)
          \bigl(1-e^{-s\log(m/\rad m)}\bigr)
\]
is nonnegative and nondecreasing in $s>0$. Therefore its value at
$w_m$ is at most its value at $w$, whose mean is $o(w)$.
We conclude that $\E_p B_{t/\log(p+2)}(p+2)\sim2C_2w$.

It remains to remove the prime-power endpoints from this mean.
The upper sieve gives $|\{p\in\mathcal P_X:p+2\in\PP\}|\ll X/L^2$
(also an instance of Lemma~\ref{lem:sieve} with all $g_\ell=0$).
Proper prime powers in $(X,2X+2]$ number $O(\sqrt X\log X)$.
Since the summands are at most one, their total is
$O(X/L^2+\sqrt X\log X)=o(Xt/L^2)$.
Finally $|\mathcal P_X|\sim X/L$ converts the mean into
\eqref{eq:interior}.
\end{proof}

\section{Mellin representation and the critical window}
For $m$ with at least two distinct prime divisors put
\begin{equation}\label{eq:signed-B}
       h_m(t)=\mu(\rad m) B_{t/\log m}(m).
\end{equation}
The sign of $h_m(t)$ is independent of $t>0$.

\begin{lemma}\label{lem:mellin}
For such an $m$ and real $0<\alpha<2$, $\alpha\ne1$,
\begin{equation}\label{eq:mellin}
 Q_\alpha(m)=\frac1{\Gamma(-\alpha)}
                \int_0^\infty h_m(t)\,\frac{dt}{t^{1+\alpha}}.
\end{equation}
Consequently
\[
 |Q_\alpha(m)|=\frac1{|\Gamma(-\alpha)|}
                \int_0^\infty B_{t/\log m}(m)\,\frac{dt}{t^{1+\alpha}}.
\]
At $\alpha=1$, $Q_1(m)=0$.
\end{lemma}

\begin{proof}
Use the usual Mellin integral for $x^\alpha$ on $x\ge0$, subtracting
$1$ from $e^{-tx}$ when $0<\alpha<1$, and subtracting $1-tx$ when
$1<\alpha<2$. The constant term disappears after summing $\mu(d)$.
The linear term also disappears because
\[
 \sum_{d\mid m}\mu(d)\log(m/d)=\Lambda(m)=0.
\]
The remaining exponential sum factors exactly as
\begin{align*}
 \sum_{d\mid m}\mu(d)e^{-t\log(m/d)/\log m}
 &=m^{-t/\log m}\prod_{\ell\mid m}(1-\ell^{t/\log m})\\
 &=\mu(\rad m)\left(\frac{\rad m}{m}\right)^{t/\log m}
                   \prod_{\ell\mid m}(1-\ell^{-t/\log m}).
\end{align*}
For this fixed $m$ the last expression is $O_m(t^2)$ at zero and
bounded at infinity. The integral converges absolutely throughout
$0<\alpha<2$. Its fixed sign proves the absolute-value formula.
The assertion at one is the displayed von Mangoldt identity.
\end{proof}

\begin{proof}[Proof of the first two limits of Theorem~\ref{th:critical}]
For the non-prime-power endpoints, Tonelli's theorem and
Lemma~\ref{lem:mellin} identify the normalized error with
\begin{equation}\label{eq:integrated-error}
 \frac1{|\Gamma(-\alpha)|}
         \int_0^\infty A_X(t)t^{-1-\alpha}\,dt.
\end{equation}
If $\alpha=1\pm c/H$, then
$H/|\Gamma(-\alpha)|\to c$.
By \eqref{eq:envelope}, the portions $0<t<1$ and $t>L$ in
\eqref{eq:integrated-error} are $O_c(H^{-1})$.
For example their integrals before the gamma factor are bounded by
$C/(2-\alpha)$ and $CL^{1-\alpha}/\alpha$, respectively.

In $1\le t\le L$, put $t=L^v$. The middle contribution is
\[
 \frac{H}{|\Gamma(-\alpha)|}
       \int_0^1\frac{A_X(L^v)}{L^v}
                   e^{-(\alpha-1)Hv}\,dv.
\]
The quotient is uniformly bounded and tends to one for each
$0<v<1$ by Proposition~\ref{prop:interior}. Dominated convergence
gives $c\int_0^1e^{-cv}\,dv=1-e^{-c}$ above one and
$c\int_0^1e^{cv}\,dv=e^c-1$ below one.

For $m=\ell^a$ with $a\ge2$, the exact formula is
$Q_\alpha(m)=1-(1-1/a)^\alpha\in[0,1]$.
Their $O(\sqrt X\log X)$ occurrences contribute $o(1)$ in
\eqref{eq:error}, completing the proof with all composite endpoints.
\end{proof}

\section{Fixed binomial approximants}
\begin{proof}[Proof of the final assertion of Theorem~\ref{th:critical}]
For every fixed $k$ and sufficiently large $X$, all orders in
\eqref{eq:Rk} lie in $(1,2)$. Their prime values are one, and
$\sum_{j=1}^k(-1)^{j+1}\binom{k}{j}=1$.
On a non-prime-power $m$ the Mellin formula and $t=L^v$ give the
middle contribution
\[
 \mu(\rad m)\int_0^1
       \frac{B_{L^v/\log m}(m)}{L^v}K_{k,H}(v)\,dv,
 \qquad
 K_{k,H}(v)=H\sum_{j=1}^k
       \frac{(-1)^{j+1}\binom{k}{j}e^{-jv}}
            {\Gamma(-1-j/H)}.
\]
Uniformly for $v\in[0,1]$,
\begin{equation}\label{eq:binomial-kernel}
 K_{k,H}(v)\longrightarrow
       k e^{-v}(1-e^{-v})^{k-1}=:K_k(v)\ge0.
\end{equation}
The two end intervals of each Mellin integral contribute $O_k(H^{-1})$
in total normalized absolute value, by the estimates already proved.
For the middle interval, replacing $K_{k,H}$ by $K_k$ changes the
normalized sum of absolute values by at most
\[
 \|K_{k,H}-K_k\|_\infty
            \int_0^1\frac{A_X(L^v)}{L^v}\,dv=o_k(1).
\]
This uses $||z|-|z'||\le|z-z'|$ for each $m$. After replacement the
integrand is nonnegative apart from the fixed sign $\mu(\rad m)$,
so its absolute value may be taken inside the integral. Dominated
convergence now gives
\[
 \int_0^1K_k(v)\,dv=(1-e^{-1})^k.
\]
Proper prime powers contribute $o_k(1)$, since $|R_{k,X}(\ell^a)|\le2^k-1$.
\end{proof}

\section{Signed limits and the role of the number of prime factors}

The same envelope also describes the signed information that the
absolute errors leave undetermined. Extend \eqref{eq:Q} to complex
$\alpha$ with $\Re\alpha>0$ using real logarithms of positive bases
and retaining the convention $0^\alpha=0$.

\begin{corollary}
Along every sequence $X\to\infty$ there is a subsequence, a number
$\tau\ge0$ and a real $b\in L^\infty(0,1)$ with $|b|\le1$ almost
everywhere such that, locally uniformly for $z\in\mathbb C$,
\[
\Mcal(Q_{1+z/H})\longrightarrow
\tau+ z\int_0^1 e^{-zv}b(v)\,dv.
\]
Here $\tau$ is the subsequential limit of the normalized twin count.
\end{corollary}
\begin{proof}
The upper sieve bounds $\tau_X$ uniformly. Define
\[
 b_X(v)=\frac{L^2}{2C_2X L^v}
 \sum_{p\in\mathcal C_X} h_{p+2}(L^v).
\]
Its absolute value is bounded by $A_X(L^v)/L^v$, uniformly bounded
on $0<v<1$ and tending pointwise to one. Pass to a subsequence with
$\tau_X\to\tau$ and $b_X\to b$ weak-star in $L^\infty(0,1)$;
separability of $L^1(0,1)$ gives sequential compactness on bounded
sets. Testing against any integrable function and using the envelope
gives $|b|\le1$ almost everywhere. For a fixed $m$ the integral in
\eqref{eq:mellin} is analytic in $0<\Re\alpha<2$, with the reciprocal
gamma factor supplying the value at one. Uniformly on compact
$z$-sets the two tails contribute $O(1/H)$ by
\eqref{eq:envelope}, and $H/\Gamma(-1-z/H)\to z$.
The middle term is the stated Laplace integral with $b_X$ in place
of $b$. Convergence is uniform on compact sets since
$z\mapsto e^{-zv}$ has compact image there in $L^1(dv)$.
The proper prime powers give $o(1)$ uniformly on these sets.
\end{proof}

\begin{corollary}
Suppose $K_X\ge2$ and $K_X=o(\log\log X)$. For either order
$\alpha=1\pm c/H$, $c>0$ fixed, the part of $\err(Q_\alpha)$ with
$2\le\omega(p+2)\le K_X$ tends to zero.
\end{corollary}
\begin{proof}
Let $A_X^{(\le K_X)}$ be \eqref{eq:AX} restricted to these endpoints.
Fix $0<v<1$, put $t=L^v$ and $w=t/L$. When $\omega(m)\le K_X$,
\[
 B_{t/\log m}(m)
 \le 2^{K_X}\!\!\!\!\prod_{\substack{\ell\mid m\\2<\ell\le X^{1/100}}}
 \frac{1-\ell^{-w}}2.
\]
Use \eqref{eq:random-upper} with these halved local weights.
By \eqref{eq:harmonic-envelope},
$A_X^{(\le K_X)}(t)\ll2^{K_X}(t^{1/2}+L^2X^{-1/2})$.
Since $2^{K_X}=L^{o(1)}$, the restricted envelope divided by $L^v$
tends to zero. Dominated convergence with the full
\eqref{eq:envelope}, including the two tails as before, proves the
claim for both signs. The proof allows repeated prime bases.
\end{proof}

There is also a precise limitation of fixed positive integer moments.
For $r$ positive logarithmic shares $x_i$ with sum one, let
\[
 q_k(x)=\sum_{S\subseteq\{1,\ldots,r\}}(-1)^{|S|}
             \left(1-\sum_{i\in S}x_i\right)^k .
\]
For integer $k\ge1$, inclusion--exclusion makes this the probability
of observing all $r$ categories in $k$ independent draws with
probabilities $x_i$. Write
\[
 a_{k,r}=\frac1{r!}\int_{\sum x_i=1,\ x_i>0}
       q_k(x)\frac{dx_1\cdots dx_{r-1}}{x_1\cdots x_r};
 \qquad a_{k,1}=1.
\]
Expanding the multinomial over compositions $j_1+\cdots+j_r=k$
with each $j_i\ge1$ and applying the simplex beta integral gives
$a_{k,r}=(k/r!)\sum_{j_1+\cdots+j_r=k}\prod_i 1/j_i$, with
positive parts understood. In particular $a_{k,r}=0$ for $r>k$,
and coefficient extraction gives
\[
 \sum_{r=1}^k a_{k,r}z^r
    = k[u^k]\exp\!\left(z\sum_{j\ge1}u^j/j\right)
    = \frac{z(z+1)\cdots(z+k-1)}{(k-1)!}.
\]
At $z=1,-1$ this shows that multiplying the angular mass at $r$
by $1+(-1)^r$ preserves its total $k$ for every $k\ge2$ while
removing all mass at $r=1$. This is an identity for the simplex
comparison measure. The preceding corollary supplies a separate
statement on the actual shifted primes.

\section{Scope, attribution and disclosure}
The upper envelope uses a fixed dimension-two upper sieve; its constants
are uniform in the local retention probabilities. Ford's total-variation
theorem supplies the interior limit for a bounded valuation statistic at
the subpower cutoff $y=\exp(L/\sqrt t)$. The gamma pole and a change of
variables then select the window $\alpha-1\asymp1/\log\log X$.
The proof does not evaluate a new signed prime correlation.

The explicit absolute-error laws and finite binomial combinations are the
proposed contribution. Their underlying distribution estimates are credited
to the cited sources. A focused source comparison found no exact stated
match for these moving-order constants, but this draft does not establish
global priority or publication acceptance. Fixed integer-moment and
probabilistic variants are not offered as separate papers.

This note was developed from an unpublished research record. ChatGPT and
AI review agents assisted with extraction, mathematical reconstruction,
literature comparison and writing. Any AI checks described in the associated
review record are internal, not external peer review or formal verification.
DEEPAK KUMAR (\href{mailto:admin@cosmosgeek.org}{admin@cosmosgeek.org}) is
the research coordinator and contact. The original derivation and subsequent checks are preserved in the \href{https://github.com/cosmosgeek0/twin-prime-research-archive}{project archive}.

\begin{thebibliography}{9}\small
\bibitem{ford-model}
K. Ford, \emph{Poisson approximation of prime divisors of shifted primes},
arXiv:2408.03803v4 (2026), Theorem 1.
\href{https://arxiv.org/html/2408.03803v4}{Readable primary text}.

\bibitem{ford-asymptotic}
K. Ford, \emph{On Bombieri's asymptotic sieve},
arXiv:math/0401215.
\href{https://arxiv.org/abs/math/0401215}{Primary preprint}.

\bibitem{kouk}
D. Koukoulopoulos, \emph{The Distribution of Prime Numbers},
American Mathematical Society, 2019, Chapters 18--19.
\href{https://dms.umontreal.ca/~koukoulo/documents/publications/primes.pdf}{Author-hosted text}.
\end{thebibliography}
\end{document}
